% 原书物理页=印刷页53--54；2.2 全节.（2.2.3 末尾命题的证明最后一行及其后一段实际印于原书第 55 页顶部、2.3 节标题之前，属本节内容，一并译出附于末尾.）
\section{Killing 型与 Casimir 算子}\label{sec:lie-2-2}

\begin{definition}\label{def:lie-2-2-1}
设 $\mathfrak g$ 是一个 $\mathbb K$-Lie 代数. $\mathfrak g$ 上一个 $\mathbb K$-双线性型 $B$ 称为\newterm{lie-invariant-bilinear-form}{不变双线性型}{invariant bilinear form}，如果 $B(\operatorname{ad}_Z(X),Y)+B(X,\operatorname{ad}_Z(Y))=0$ 对所有 $X,Y,Z\in\mathfrak g$ 成立.
\end{definition}

\begin{example*}
不变双线性型的范式是 $\mathfrak{gl}_n(\mathbb K)$ 上的型
\[
(X,Y)\mapsto\operatorname{tr}(XY).
\]
这给出了构造许多其它此类型的方法：对任何表示 $r:\mathfrak g\longrightarrow\mathfrak{gl}(V)$，$\mathfrak g$ 上的双线性型 $B_r(X,Y):=\operatorname{tr}(r(X)r(Y))$ 是不变的.
\end{example*}

\begin{proposition*}
若 $B$ 是 $\mathfrak g$ 上的不变双线性型，而 $\mathfrak h\subset\mathfrak g$ 是 $\mathfrak g$ 的一个理想，则它的\newterm{lie-orthogonal-complement}{正交补}{orthogonal} $\mathfrak h^{\perp}=\{X\in\mathfrak g,\ \forall Y\in\mathfrak h,\ B(X,Y)=0\}$ 也是 $\mathfrak g$ 的理想. 特别地，$\mathfrak g^{\perp}$ 是 $\mathfrak g$ 的理想.
\end{proposition*}

\noindent\textbf{Killing 型.}\quad $\mathfrak g$ 上的\newterm{lie-killing-form}{Killing 型}{Killing form} $B_{\mathrm{ad}}$ 是与表示 $r=\operatorname{ad}$ 相配（如上）的不变双线性型. 于是 $B_{\mathrm{ad}}(X,Y)=\operatorname{tr}(\operatorname{ad}_X\operatorname{ad}_Y)$.

\begin{lemma}\label{lemma:lie-2-2-2}
若 $\mathfrak h$ 是 $\mathfrak g$ 的一个理想\footnote{注意，若 $\mathfrak h$ 只是 $\mathfrak g$ 的一个 Lie 子代数，结论并不成立.}，则 $\mathfrak h$ 的 Killing 型是 $\mathfrak g$ 的 Killing 型在 $\mathfrak h$ 上的限制.
\end{lemma}

\begin{proof}
对 $X,Y\in\mathfrak h$，有 $\operatorname{ad}_X\operatorname{ad}_Y(\mathfrak g)\subset\mathfrak h$. 故 $\operatorname{tr}(\operatorname{ad}_X\operatorname{ad}_Y|_{\mathfrak g})=\operatorname{tr}(\operatorname{ad}_X\operatorname{ad}_Y|_{\mathfrak h})$.
\end{proof}

\begin{numpara}[非退化不变双线性型]\label{np:lie-2-2-3}
现在设已给定 $\mathfrak g$ 上一个非退化的双线性型 $B$. 我们给它配上泛包络代数中的一个元素 $C_B\in\mathcal U\mathfrak g$，称为 Casimir 算子，定义为
\[
C_B=\sum_{i=1}^{n}X_iX_i^{*},
\]
其中 $X_1,\cdots,X_n$ 表示 $\mathfrak g$ 的一组基，而 $X_1^{*},\cdots,X_n^{*}$ 表示关于非退化型 $B$ 的\newterm{lie-dual-basis}{对偶基}{dual basis}.
\end{numpara}

\begin{lemma*}
此定义不依赖基的选取.
\end{lemma*}

\begin{proof}
设 $Y_1,\cdots,Y_n$ 是 $\mathfrak g$ 的另一组基. 我们有 $Y_i=\sum_k B(Y_i,X_k^{*})X_k$ 以及 $Y_i^{*}=\sum_l B(Y_i^{*},X_l)X_l^{*}$. 于是
\begin{align*}
\sum_i Y_iY_i^{*}
&=\sum_{k,l}\Bigl(\sum_i B(Y_i,X_k^{*})B(Y_i^{*},X_l)\Bigr)X_kX_l^{*}\\
&=\sum_{k,l}B\Bigl(\sum_i Y_iB(Y_i^{*},X_l),\,X_k^{*}\Bigr)X_kX_l^{*}
=\sum_{k,l}B(X_l,X_k^{*})X_kX_l^{*}=C_B.
\end{align*}
\end{proof}

\begin{example*}
取 $\mathfrak g=\mathfrak{sl}_2(\mathbb K)$，$B=B_{\mathrm{std}}$ 为由 $(X,Y)\mapsto\operatorname{tr}(XY)$ 定义的标准双线性型. 容易验证它非退化，且基 $(E,H,F)$ 的对偶基是 $(F,\tfrac12H,E)$. 我们重新得到了前面引入的算子 $C=EF+\tfrac12H^2+FE$.
\end{example*}

\begin{exercise*}
\begin{enumerate}[label=\roman*)]
\item 验证 $(X,Y)\mapsto\operatorname{tr}(XY)$ 在 $\mathfrak{gl}_n(\mathbb K)$ 上以及在经典 Lie 代数（$\mathfrak{so}$、$\mathfrak{sl}$、$\mathfrak{sp}$）上非退化.
\item （$\mathbb K=\mathbb R$ 或 $\mathbb C$.）证明：对 $\mathfrak{gl}_n(\mathbb R)$ 的任何在转置下稳定的 Lie 子代数 $\mathfrak g$，双线性型 $(X,Y)\mapsto\operatorname{tr}(XY)$ 非退化.
\end{enumerate}
\end{exercise*}

我们已经看到，$\mathfrak{sl}_2(\mathbb R)$ 的 Casimir 算子的妙处在于它属于 $\mathcal U\mathfrak{sl}_2(\mathbb R)$ 的中心. 更一般地：

\begin{proposition*}
算子 $C_B$ 属于 $\mathcal U\mathfrak g$ 的中心.
\end{proposition*}

\begin{proof}
只需证明它与任何 $X\in\mathfrak g$ 交换. 而我们有
\begin{align*}
XC_B-C_BX
&=\sum_i\bigl(XX_iX_i^{*}-X_iX_i^{*}X\bigr)
=\sum_i\bigl([X,X_i]X_i^{*}-X_i[X_i^{*},X]\bigr)\\
&=\sum_i\sum_k\bigl(B([X,X_i],X_k^{*})X_kX_i^{*}-B(X_k,[X_i^{*},X])X_iX_k^{*}\bigr)\\
&=\sum_{i,k}\bigl(B([X,X_k],X_i^{*})+B(X_k,[X,X_i^{*}])\bigr)X_iX_k^{*}=0
\end{align*}
因为双线性型 $B$ 不变.
\end{proof}

与 $\mathfrak{sl}_2$ 的情形一样，一个有趣的推论是：在 $\mathfrak g$ 的任何表示 $(V,r)$ 上，元素 $C$ 诱导一个算子 $r(C)=\sum_i r(X_i)r(X_i^{*})$，它与 $\mathfrak g$ 的作用交换. 特别地，若 $V$ 不可约且 $\mathbb K$ 代数闭（例如 $\mathbb K=\mathbb C$），则 $r(C)$ 是一个位似.
